% Documento de cálculos que entregó el usuario (planta ANTERIOR). Sirve de plantilla de formato.
\documentclass[11pt]{article}
\usepackage[utf8]{inputenc}
\usepackage[T1]{fontenc}
\usepackage[spanish,es-noshorthands,es-nodecimaldot]{babel}
\usepackage{amsmath,amssymb}
\usepackage[margin=2cm]{geometry}
\setcounter{MaxMatrixCols}{12}
\allowdisplaybreaks
\begin{document}
\section*{Cálculos -- Laboratorio 4}

\subsection*{1. Planta}
\begin{gather*}
G(s)=\frac{X(s)}{U(s)}=\frac{270.02}{s^2+29.13s+57.45}\\
s^2+29.13s+57.45=0 \;\Rightarrow\; s_1=-2.128,\quad s_2=-27.00\\
K_{dc}=\frac{270.02}{57.45}=4.700~\text{mm/V},\qquad t_{s,la}(2\%)=1.877~\text{s}\\
x_1=x,\quad x_2=\dot{x}\\
A=\begin{bmatrix}0&1\\-57.45&-29.13\end{bmatrix},\quad
B=\begin{bmatrix}0\\270.02\end{bmatrix},\quad
C=\begin{bmatrix}1&0\end{bmatrix},\quad D=0
\end{gather*}
\subsection*{2. Controlabilidad y observabilidad}
\begin{gather*}
S=[B\;\;AB]=\begin{bmatrix}0 & 270.02 \\ 270.02 & -7865.7\end{bmatrix},\quad \det S=-72911\neq0\\
\mathcal{O}=\begin{bmatrix}C\\CA\end{bmatrix}=\begin{bmatrix}1 & 0 \\ 0 & 1\end{bmatrix},\quad \det\mathcal{O}=1\neq0
\end{gather*}
\subsection*{3. Especificaciones}
\begin{gather*}
t_{sd}=1~\text{s},\qquad \zeta_d=0.8,\qquad \omega_{nd}=\frac{4}{\zeta_d t_{sd}}=\frac{4}{0.8(1)}=5~\text{rad/s}\\
s^2+2(0.8)(5)s+5^2=s^2+8s+25 \;\Rightarrow\; s=-4\pm3j\\
p_{\text{lejano}}=10\,\zeta_d\omega_{nd}=10(0.8)(5)=40
\end{gather*}
\subsection*{4. Controlador Escalón ($m=1$)}
\[\dot\rho=r-y,\qquad u=-K\mathbf{x}+K_i\rho,\qquad \mathbf{z}=[x_1\;\;x_2\;\;\rho]^T\]
\begin{gather*}
\hat{A}=\begin{bmatrix}0 & 1 & 0 \\ -57.45 & -29.13 & 0 \\ -1 & 0 & 0\end{bmatrix},\quad \hat{B}=\begin{bmatrix}0 \\ 270.02 \\ 0\end{bmatrix}\\
\hat{S}=\begin{bmatrix}0 & 270.02 & -7865.7 \\ 270.02 & -7865.7 & 213614.7 \\ 0 & 0 & -270.02\end{bmatrix}\\
\det\hat{S}=1.969\times10^{7}\neq0\\
P_d(s)=(s^2+8s+25)(s+40)=s^{3} + 48s^{2} + 345s + 1000
\end{gather*}
\paragraph{Método 1: asignación de polos}
\begin{gather*}
A^*=\begin{bmatrix}A-BK&BK_i\\-C&0\end{bmatrix}=\begin{bmatrix}0&1&0\\-57.45-270.02K_1&-29.13-270.02K_2&270.02K_i\\-1&0&0\end{bmatrix}\\
sI-A^*=\begin{bmatrix}s&-1&0\\57.45+270.02K_1&s+29.13+270.02K_2&-270.02K_i\\1&0&s\end{bmatrix}\\
\det(sI-A^*)=s^3+(29.13+270.02K_2)s^2+(57.45+270.02K_1)s+270.02K_i
\end{gather*}
\begin{align*}
29.13+270.02K_2=48 &\Rightarrow K_2=0.06988\\
57.45+270.02K_1=345 &\Rightarrow K_1=1.0649\\
270.02K_i=1000 &\Rightarrow K_i=3.7034
\end{align*}
\[\hat{K}=\begin{bmatrix}1.0649 & 0.069884 & -3.7034\end{bmatrix}\]
\paragraph{Método 2: Ackermann}
\begin{gather*}
\phi(\hat{A})=\hat{A}^{3} + 48\hat{A}^{2} + 345\hat{A} + 1000I\\
\phi(\hat{A})=\begin{bmatrix}-84.082 & -262.13 & 0 \\ 15060 & 7551.9 & 0 \\ -287.55 & -18.87 & 1000\end{bmatrix}\\
[0\;\cdots\;0\;\;1]\,\hat{S}^{-1}=\begin{bmatrix}0 & 0 & -0.0037034\end{bmatrix}\\
\hat{K}=[0\;\cdots\;0\;\;1]\,\hat{S}^{-1}\phi(\hat{A})=\begin{bmatrix}1.0649 & 0.069884 & -3.7034\end{bmatrix}
\end{gather*}
\paragraph{Método 3: matriz de transformación}
\begin{gather*}
\det(sI-\hat{A})=s^{3} + 29.13s^{2} + 57.45s\\
\hat{M}=\begin{bmatrix}57.45 & 29.13 & 1 \\ 29.13 & 1 & 0 \\ 1 & 0 & 0\end{bmatrix}\\
T=\hat{S}\hat{M}=\begin{bmatrix}0 & 270.02 & 0 \\ 0 & 0 & 270.02 \\ -270.02 & 0 & 0\end{bmatrix}\\
T^{-1}=\begin{bmatrix}0 & 0 & -0.0037034 \\ 0.0037034 & 0 & 0 \\ 0 & 0.0037034 & 0\end{bmatrix}\\
[\alpha_3-a_3\;\cdots\;\alpha_1-a_1]=\begin{bmatrix}1000 & 287.55 & 18.87\end{bmatrix}\\
\hat{K}=\begin{bmatrix}1000 & 287.55 & 18.87\end{bmatrix}\,T^{-1}\\
\hat{K}=\begin{bmatrix}1.0649 & 0.069884 & -3.7034\end{bmatrix}\\
\text{eig}(\hat{A}-\hat{B}\hat{K})=\{-40, -4\pm3j\}
\end{gather*}
\subsection*{5. Controlador Rampa ($m=2$)}
\[\ddot\rho=r-y,\qquad u=-K\mathbf{x}+K_i\dot\rho+K_{i2}\rho,\qquad \mathbf{z}=[x_1\;\;x_2\;\;\rho\;\;\dot\rho]^T\]
\begin{gather*}
\hat{A}=\begin{bmatrix}0 & 1 & 0 & 0 \\ -57.45 & -29.13 & 0 & 0 \\ 0 & 0 & 0 & 1 \\ -1 & 0 & 0 & 0\end{bmatrix},\quad \hat{B}=\begin{bmatrix}0 \\ 270.02 \\ 0 \\ 0\end{bmatrix}\\
\hat{S}=\begin{bmatrix}0 & 270.02 & -7865.7 & 213614.7 \\ 270.02 & -7865.7 & 213614.7 & -5770712 \\ 0 & 0 & 0 & -270.02 \\ 0 & 0 & -270.02 & 7865.7\end{bmatrix}\\
\det\hat{S}=5.316\times10^{9}\neq0\\
P_d(s)=(s^2+8s+25)(s+40)^2=s^{4} + 88s^{3} + 2265s^{2} + 14800s + 40000
\end{gather*}
\paragraph{Método 1: asignación de polos}
\begin{gather*}
A^*=\begin{bmatrix}A-BK&BK_{i2}&BK_i\\0&0&1\\-C&0&0\end{bmatrix}=\begin{bmatrix}0&1&0&0\\-57.45-270.02K_1&-29.13-270.02K_2&270.02K_{i2}&270.02K_i\\0&0&0&1\\-1&0&0&0\end{bmatrix}\\
\det(sI-A^*)=s^4+(29.13+270.02K_2)s^3+(57.45+270.02K_1)s^2+270.02K_i\,s+270.02K_{i2}
\end{gather*}
\begin{align*}
29.13+270.02K_2=88 &\Rightarrow K_2=0.21802\\
57.45+270.02K_1=2265 &\Rightarrow K_1=8.1755\\
270.02K_i=14800 &\Rightarrow K_i=54.8108\\
270.02K_{i2}=40000 &\Rightarrow K_{i2}=148.1372
\end{align*}
\[\hat{K}=\begin{bmatrix}8.1755 & 0.21802 & -148.14 & -54.811\end{bmatrix}\]
\paragraph{Método 2: Ackermann}
\begin{gather*}
\phi(\hat{A})=\hat{A}^{4} + 88\hat{A}^{3} + 2265\hat{A}^{2} + 14800\hat{A} + 40000I\\
\phi(\hat{A})=\begin{bmatrix}11696 & -2933.5 & 0 & 0 \\ 168527.8 & 97148 & 0 & 0 \\ -2207.6 & -58.87 & 40000 & 14800 \\ -11418 & -492.67 & 0 & 40000\end{bmatrix}\\
[0\;\cdots\;0\;\;1]\,\hat{S}^{-1}=\begin{bmatrix}0 & 0 & -0.0037034 & 0\end{bmatrix}\\
\hat{K}=[0\;\cdots\;0\;\;1]\,\hat{S}^{-1}\phi(\hat{A})=\begin{bmatrix}8.1755 & 0.21802 & -148.14 & -54.811\end{bmatrix}
\end{gather*}
\paragraph{Método 3: matriz de transformación}
\begin{gather*}
\det(sI-\hat{A})=s^{4} + 29.13s^{3} + 57.45s^{2}\\
\hat{M}=\begin{bmatrix}0 & 57.45 & 29.13 & 1 \\ 57.45 & 29.13 & 1 & 0 \\ 29.13 & 1 & 0 & 0 \\ 1 & 0 & 0 & 0\end{bmatrix}\\
T=\hat{S}\hat{M}=\begin{bmatrix}0 & 0 & 270.02 & 0 \\ 0 & 0 & 0 & 270.02 \\ -270.02 & 0 & 0 & 0 \\ 0 & -270.02 & 0 & 0\end{bmatrix}\\
T^{-1}=\begin{bmatrix}0 & 0 & -0.0037034 & 0 \\ 0 & 0 & 0 & -0.0037034 \\ 0.0037034 & 0 & 0 & 0 \\ 0 & 0.0037034 & 0 & 0\end{bmatrix}\\
[\alpha_4-a_4\;\cdots\;\alpha_1-a_1]=\begin{bmatrix}40000 & 14800 & 2207.6 & 58.87\end{bmatrix}\\
\hat{K}=\begin{bmatrix}40000 & 14800 & 2207.6 & 58.87\end{bmatrix}\,T^{-1}\\
\hat{K}=\begin{bmatrix}8.1755 & 0.21802 & -148.14 & -54.811\end{bmatrix}\\
\text{eig}(\hat{A}-\hat{B}\hat{K})=\{-40, -4\pm3j\}
\end{gather*}
\subsection*{6. Controlador Parábola ($m=3$)}
\[\dddot\rho=r-y,\qquad u=-K\mathbf{x}+K_i\ddot\rho+K_{i2}\dot\rho+K_{i3}\rho,\qquad \mathbf{z}=[x_1\;\;x_2\;\;\rho\;\;\dot\rho\;\;\ddot\rho]^T\]
\begin{gather*}
\hat{A}=\begin{bmatrix}0 & 1 & 0 & 0 & 0 \\ -57.45 & -29.13 & 0 & 0 & 0 \\ 0 & 0 & 0 & 1 & 0 \\ 0 & 0 & 0 & 0 & 1 \\ -1 & 0 & 0 & 0 & 0\end{bmatrix},\quad \hat{B}=\begin{bmatrix}0 \\ 270.02 \\ 0 \\ 0 \\ 0\end{bmatrix}\\
\hat{S}=\begin{bmatrix}0 & 270.02 & -7865.7 & 213614.7 & -5770712 \\ 270.02 & -7865.7 & 213614.7 & -5770712 & 1.558\times10^{8} \\ 0 & 0 & 0 & 0 & -270.02 \\ 0 & 0 & 0 & -270.02 & 7865.7 \\ 0 & 0 & -270.02 & 7865.7 & -213614.7\end{bmatrix}\\
\det\hat{S}=-1.435\times10^{12}\neq0\\
P_d(s)=(s^2+8s+25)(s+40)^3=s^{5} + 128s^{4} + 5785s^{3} + 105400s^{2} + 632000s + 1600000
\end{gather*}
\paragraph{Método 1: asignación de polos}
\begin{gather*}
A^*=\begin{bmatrix}A-BK&BK_{i3}&BK_{i2}&BK_i\\0&0&1&0\\0&0&0&1\\-C&0&0&0\end{bmatrix}\\
\det(sI-A^*)=s^5+(29.13+270.02K_2)s^4+(57.45+270.02K_1)s^3\\+270.02K_i\,s^2+270.02K_{i2}\,s+270.02K_{i3}
\end{gather*}
\begin{align*}
29.13+270.02K_2=128 &\Rightarrow K_2=0.36616\\
57.45+270.02K_1=5785 &\Rightarrow K_1=21.2116\\
270.02K_i=105400 &\Rightarrow K_i=390.3415\\
270.02K_{i2}=632000 &\Rightarrow K_{i2}=2340.5674\\
270.02K_{i3}=1600000 &\Rightarrow K_{i3}=5925.487
\end{align*}
\[\hat{K}=\begin{bmatrix}21.212 & 0.36616 & -5925.5 & -2340.6 & -390.34\end{bmatrix}\]
\paragraph{Método 2: Ackermann}
\begin{gather*}
\phi(\hat{A})=\hat{A}^{5} + 128\hat{A}^{4} + 5785\hat{A}^{3} + 105400\hat{A}^{2} + 632000\hat{A} + 1600000I\\
\phi(\hat{A})=\begin{bmatrix}636379 & -20191 & 0 & 0 & 0 \\ 1159945 & 1224529 & 0 & 0 & 0 \\ -5727.6 & -98.87 & 1600000 & 632000 & 105400 \\ -99720 & -2847.5 & 0 & 1600000 & 632000 \\ -468413 & -16773 & 0 & 0 & 1600000\end{bmatrix}\\
[0\;\cdots\;0\;\;1]\,\hat{S}^{-1}=\begin{bmatrix}0 & 0 & -0.0037034 & 0 & 0\end{bmatrix}\\
\hat{K}=[0\;\cdots\;0\;\;1]\,\hat{S}^{-1}\phi(\hat{A})=\begin{bmatrix}21.212 & 0.36616 & -5925.5 & -2340.6 & -390.34\end{bmatrix}
\end{gather*}
\paragraph{Método 3: matriz de transformación}
\begin{gather*}
\det(sI-\hat{A})=s^{5} + 29.13s^{4} + 57.45s^{3}\\
\hat{M}=\begin{bmatrix}0 & 0 & 57.45 & 29.13 & 1 \\ 0 & 57.45 & 29.13 & 1 & 0 \\ 57.45 & 29.13 & 1 & 0 & 0 \\ 29.13 & 1 & 0 & 0 & 0 \\ 1 & 0 & 0 & 0 & 0\end{bmatrix}\\
T=\hat{S}\hat{M}=\begin{bmatrix}0 & 0 & 0 & 270.02 & 0 \\ 0 & 0 & 0 & 0 & 270.02 \\ -270.02 & 0 & 0 & 0 & 0 \\ 0 & -270.02 & 0 & 0 & 0 \\ 0 & 0 & -270.02 & 0 & 0\end{bmatrix}\\
T^{-1}=\begin{bmatrix}0 & 0 & -0.0037034 & 0 & 0 \\ 0 & 0 & 0 & -0.0037034 & 0 \\ 0 & 0 & 0 & 0 & -0.0037034 \\ 0.0037034 & 0 & 0 & 0 & 0 \\ 0 & 0.0037034 & 0 & 0 & 0\end{bmatrix}\\
[\alpha_5-a_5\;\cdots\;\alpha_1-a_1]=\begin{bmatrix}1600000 & 632000 & 105400 & 5727.6 & 98.87\end{bmatrix}\\
\hat{K}=\begin{bmatrix}1600000 & 632000 & 105400 & 5727.6 & 98.87\end{bmatrix}\,T^{-1}\\
\hat{K}=\begin{bmatrix}21.212 & 0.36616 & -5925.5 & -2340.6 & -390.34\end{bmatrix}\\
\text{eig}(\hat{A}-\hat{B}\hat{K})=\{-40, -4\pm3j\}
\end{gather*}
\subsection*{7. Observador}
\begin{gather*}
t_{s,o}=\frac{t_{sd}}{10}=0.1~\text{s},\qquad \zeta_o=0.8,\qquad \omega_{no}=\frac{4}{0.8(0.1)}=50~\text{rad/s}\\
P_o(s)=s^2+2(0.8)(50)s+50^2=s^2+80s+2500\;\Rightarrow\; s=-40\pm30j
\end{gather*}
\paragraph{Método 1}
\begin{gather*}
A-LC=\begin{bmatrix}-L_1&1\\-57.45-L_2&-29.13\end{bmatrix},\qquad
sI-(A-LC)=\begin{bmatrix}s+L_1&-1\\57.45+L_2&s+29.13\end{bmatrix}\\
\det=s^2+(L_1+29.13)s+(29.13L_1+57.45+L_2)
\end{gather*}
\begin{align*}
L_1+29.13=80 &\Rightarrow L_1=50.87\\
29.13(50.87)+57.45+L_2=2500 &\Rightarrow L_2=960.71
\end{align*}
\paragraph{Método 2: matriz de transformación}
\begin{gather*}
M=\begin{bmatrix}29.13 & 1 \\ 1 & 0\end{bmatrix},\quad Q=(M\mathcal{O})^{-1}=\begin{bmatrix}0 & 1 \\ 1 & -29.13\end{bmatrix}\\
L=Q\begin{bmatrix}2500-57.45\\80-29.13\end{bmatrix}=\begin{bmatrix}0 & 1 \\ 1 & -29.13\end{bmatrix}\begin{bmatrix}2442.55\\50.87\end{bmatrix}=\begin{bmatrix}50.87 \\ 960.71\end{bmatrix}
\end{gather*}
\paragraph{Método 3: Ackermann}
\begin{gather*}
\phi(A)=A^2+80A+2500I=\begin{bmatrix}2442.5 & 50.87 \\ -2922.5 & 960.71\end{bmatrix}\\
L=\phi(A)\,\mathcal{O}^{-1}\begin{bmatrix}0\\1\end{bmatrix}=\begin{bmatrix}50.87 \\ 960.71\end{bmatrix}\\
\text{eig}(A-LC)=\{-40\pm30j\}
\end{gather*}
\subsection*{8. Error en estado estable}
\begin{gather*}
\frac{E(s)}{R(s)}=\frac{s^m\left[s^2+(a_1+bK_2)s+(a_2+bK_1)\right]}{P_d(s)}\\
m=1,\; r=\dot r\,t:\quad e_{ss}=\frac{345}{1000}\,\dot r=0.345\,\dot r\;\;\Rightarrow\;\; \dot r=2~\text{mm/s}\Rightarrow e_{ss}=0.69~\text{mm}\\
m=2,\; r=\tfrac12\ddot r\,t^2:\quad e_{ss}=\frac{2265}{40000}\,\ddot r=0.0566\,\ddot r\\
m=2,\; r=\dot r\,t:\; e_{ss}=0\qquad m=3,\; r=\tfrac12\ddot r\,t^2:\; e_{ss}=0
\end{gather*}
\subsection*{9. Implementación con amplificadores operacionales}
\begin{gather*}
S_x=\frac{3.3~\text{V}}{8~\text{mm}}=0.4125~\text{V/mm},\qquad S_v=0.1~\text{V/(mm/s)},\qquad V_r=S_x\,r\\
P=S_x\,x,\qquad V=S_v\,\dot x,\qquad E=V_r-P=S_x(r-x)\\
I=-\frac{1}{R_{10}C_1}\int E\,dt=-\frac{S_x}{R_{10}C_1}\,\rho\\
u=-\frac{R_{20}}{R_{14}}P-\frac{R_{20}}{R_{13}}V-\frac{R_{20}}{R_{12}}I
=-\underbrace{\frac{R_{20}}{R_{14}}S_x}_{K_1}x-\underbrace{\frac{R_{20}}{R_{13}}S_v}_{K_2}\dot x+\underbrace{\frac{R_{20}}{R_{12}}\frac{S_x}{R_{10}C_1}}_{K_i}\rho
\end{gather*}
\begin{gather*}
R_{20}=R_{10}=100~\text{k}\Omega,\qquad C_1=1~\mu\text{F},\qquad R_{10}C_1=0.1~\text{s}\\
R_{14}=\frac{R_{20}S_x}{K_1}=\frac{(100\,\text{k})(0.4125)}{1.0649}=38.74~\text{k}\Omega\;\to\;39~\text{k}\Omega\\
R_{13}=\frac{R_{20}S_v}{K_2}=\frac{(100\,\text{k})(0.1)}{0.06988}=143.09~\text{k}\Omega\;\to\;143~\text{k}\Omega\\
R_{12}=\frac{R_{20}S_x}{K_i\,R_{10}C_1}=\frac{(100\,\text{k})(0.4125)}{(3.7034)(0.1)}=111.38~\text{k}\Omega\;\to\;111.5~\text{k}\Omega\\
K_1=\frac{100}{39}(0.4125)=1.0577,\quad K_2=\frac{100}{143}(0.1)=0.06993,\quad K_i=\frac{100}{111.5}\cdot\frac{0.4125}{0.1}=3.6996\\
\text{eig}(\hat{A}-\hat{B}\hat{K}_{\text{real}})=\{-40.07, -3.97\pm3.03j\}
\end{gather*}
\end{document}
